\documentclass[11pt,a4paper]{article}
\usepackage[T1]{fontenc}
\usepackage{lmodern,amsmath,amssymb,amsthm,mathtools,booktabs,tabularx,array}
\usepackage[a4paper,margin=27mm,top=25mm,bottom=25mm]{geometry}
\usepackage{microtype,xcolor,fancyhdr,hyperref}
\definecolor{navy}{HTML}{17364B}
\definecolor{teal}{HTML}{116E70}
\hypersetup{colorlinks=true,linkcolor=teal,citecolor=teal,urlcolor=teal,pdftitle={Arithmetic-Geometric and Galois Rigidity of Tate-Shafarevich Torsors},pdfauthor={João Víctor de Mello}}
\pagestyle{fancy}\fancyhf{}\fancyhead[L]{\small\textsc{Arithmetic--geometric rigidity}}\fancyhead[R]{\small J. V. de Mello}\fancyfoot[C]{\thepage}\setlength{\headheight}{14pt}
\newtheorem{theorem}{Theorem}[section]
\newtheorem{lemma}[theorem]{Lemma}
\newtheorem{corollary}[theorem]{Corollary}
\theoremstyle{definition}\newtheorem{definition}[theorem]{Definition}
\theoremstyle{remark}\newtheorem{remark}[theorem]{Remark}
\newcommand{\Sha}{\operatorname{\mathrm{Sha}}}
\newcommand{\ZZ}{\mathbb Z}\newcommand{\QQ}{\mathbb Q}\newcommand{\FF}{\mathbb F}
\newcommand{\Gal}{\operatorname{Gal}}\newcommand{\Hom}{\operatorname{Hom}}\newcommand{\End}{\operatorname{End}}\newcommand{\Pic}{\operatorname{Pic}}\newcommand{\NS}{\operatorname{NS}}\newcommand{\Br}{\operatorname{Br}}\newcommand{\Sel}{\operatorname{Sel}}\newcommand{\ord}{\operatorname{ord}}\newcommand{\rk}{\operatorname{rank}}\newcommand{\corank}{\operatorname{corank}}\newcommand{\length}{\operatorname{length}}\newcommand{\Sym}{\operatorname{Sym}}\newcommand{\Cl}{\operatorname{Cl}}\newcommand{\coker}{\operatorname{coker}}\newcommand{\I}{\mathfrak I}\newcommand{\D}{\mathfrak D}\newcommand{\Tdef}{\mathfrak T}
\numberwithin{equation}{section}
\setlength{\parindent}{1.1em}\setlength{\parskip}{2pt}
\begin{document}
\thispagestyle{plain}
\begin{center}
{\small\color{teal}\textsc{Arithmetic geometry \;\textbullet\; Research manuscript}}\\[10pt]
{\LARGE\bfseries\color{navy} Arithmetic--Geometric and Galois Rigidity of Tate--Shafarevich Torsors}\\[7pt]
{\large Polarization lattices, isogeny descent, Selmer division fields and integral motivic obstructions}\\[13pt]
{\large Jo\~ao V\'{\i}ctor de Mello}\\[4pt]
26 November 2026
\end{center}
\begin{abstract}
Let $E$ be a geometrically non-CM elliptic curve over a number field $K$. For an everywhere locally soluble torsor $X$ under $E^h$, let $H$ be the finite subgroup of $\Sha(E/K)$ generated by its coordinate classes. We identify the minimum Euler characteristic of a descended ample line bundle with $\#H$ and prove that the minimum degree of a dominant morphism from $X$ to an arbitrary abelian variety of dimension $h$ lies between $(\#H)^2/C$ and $(\#H)^2$, for a constant depending only on $E,K,h$. At primes with full residual matrix algebra the comparison is exact. We compute mixed Poincar\'{e} descent defects, integral cycle-class cokernels and canonical-height energy bounds, and derive bounded-degree finiteness. In a complementary Galois-theoretic direction, we realize the dual of $H$ with coefficients in $E[p^n]$ as a relative Selmer division-field Galois group. At a large-image prime the degree of this extension agrees exactly with the minimum geometric descent degree and the square of the minimum polarization Euler characteristic. We develop the infinite-level Lie-algebraic and fine Selmer interpretations and construct locally trivial Brauer classes from compatible torsor towers. The resulting identities quantify the divisible obstruction but do not establish unrestricted finiteness of $\Sha$ or a higher Kato motivic reciprocity formula.
\end{abstract}
\noindent\textbf{Keywords:} elliptic curves, Tate--Shafarevich groups, torsors, polarizations, Galois representations, Selmer groups, integral motives, Brauer groups.\\
\textbf{MSC 2020:} 11G05, 11G10, 14K05, 14F22, 11R23, 11G40.

\section{Introduction and scope}
The Birch and Swinnerton-Dyer conjecture predicts, among other assertions, the finiteness of the Tate--Shafarevich group. This article compares several integral invariants that retain the same arithmetic obstruction in different categories. The geometry is that of a fixed power of one elliptic curve, while the descent data are allowed to vary. No assumption of finiteness of the full Tate--Shafarevich group is made.

The starting point is the description of N\'{e}ron--Severi classes on products of elliptic curves by symmetric matrices, as in \cite{RS}, together with Picard descent and the Brauer--Hasse--Noether theorem. The isogeny estimate uses the integral matrix algebra generated by the Galois image, in contrast to an estimate for a fixed target alone. The division-field construction uses finite-module duality and a scalar annihilator to apply Kummer-type evaluation to all Selmer classes, not only classes of rational points. Related class-field methods are developed by Hiranouchi and Ohshita \cite{HO}. The Brauer construction uses compatible multiplication coverings; the relation between divisible Albanese classes and Galois sections is studied by Harari and Szamuely \cite{HS}. The Kato discussion uses existing structure results \cite{BCGS,Kim}, not a new higher explicit reciprocity theorem.

\subsection*{Summary of principal comparisons}
For $X$ under $E^h$ and $H$ as in the abstract, put $\mu(X)=\min_{L\text{ ample}}\chi(X,L)$. The notation in the table is defined precisely below.
\begin{center}\small
\begin{tabularx}{\textwidth}{@{}p{.39\textwidth}X@{}}
\toprule
\textbf{Result} & \textbf{Identity or estimate}\\ \midrule
Polarization and stable descent & $\mu(X\times E^t)=\#H$\\[3pt]
Quantitative isogeny rigidity & $(\#H)^2/C\leq\I(X)\leq(\#H)^2$\\[3pt]
Exact large-image rigidity & $\I(X)=(\#H)^2$ for $p$-primary $H$ at a large-image prime\\[3pt]
Canonical-height energy & $\min_{f:X\to E^h}\mathcal E_P(f)=\sum_i q_i^2$ for independent classes\\[3pt]
Mixed Poincar\'{e} descent & $M_h(\ZZ)/\mathcal P_K(X,A)\simeq H^{\oplus h}$\\[3pt]
Integral motive separation & $h(X)_{\ZZ_p}\not\simeq h(A)_{\ZZ_p}$ if $p\mid\#H$\\[3pt]
Bounded-degree finiteness & $\I(X)\leq D\ \Rightarrow\ \#H\leq\sqrt{CD}$\\[3pt]
Relative Selmer division field & $\Gal(L_{H,n}/L_{\mathrm{MW},n})\simeq\Hom(H,E[p^n])$\\[3pt]
Three-way comparison & $[L_{H,n}:L_{\mathrm{MW},n}]=\I(X)=\mu(X)^2$ at a large-image prime\\[3pt]
Divisible tower obstruction & $T_p\mathcal B(X_n)\twoheadrightarrow(\ZZ/p^n)^{h(h+1)/2}$\\
\bottomrule
\end{tabularx}
\end{center}
Here and below $q_i$ denotes an invariant factor or, when stated, the order of an independent coordinate torsor. The adjective ``large-image'' in the Selmer statements means $p\geq5$ and $\rho_{E,p}(G_K)\supseteq\mathrm{SL}_2(\FF_p)$. All unconditional conclusions are subject to the explicit hypotheses of their statements. The novelty of the combination relative to the entire literature is not established; the paper is a candidate research manuscript.

\section{Arithmetic polarization lattices}
Fix a geometrically non-CM elliptic curve $E/K$, so $\End_{\overline K}(E)=\ZZ$. Let $A=E^h$, equipped with its product principal polarization. Let $X$ be the $A$-torsor of class $\xi=(\xi_1,\ldots,\xi_h)$, with each $\xi_i\in\Sha(E/K)$. Define
\[
H=\langle\xi_1,\ldots,\xi_h\rangle,\qquad
\Lambda_\xi=\ker\bigl(\ZZ^h\longrightarrow H,\ a\longmapsto\textstyle\sum_i a_i\xi_i\bigr).
\]
All torsor classes are torsion, so $H$ is finite and $[\ZZ^h:\Lambda_\xi]=\#H$. Over $\overline K$, a choice of a point of $X$ identifies its N\'{e}ron--Severi group with $\Sym_h(\ZZ)$. Changing the point translates the variety and does not change this identification. The action of $G_K$ on the N\'{e}ron--Severi group is trivial. Define $\NS_K(X)$ to be the image of $\Pic(X)$ in this group.

\begin{theorem}[Polarization descent]\label{thm:pol}
Under the preceding hypotheses,
\begin{equation}\label{eq:pol}
\NS_K(X)=\{B\in\Sym_h(\ZZ): B\ZZ^h\subseteq\Lambda_\xi\}.
\end{equation}
For a line bundle with class $B$, $\chi(X,L)=\det B$; ampleness is equivalent to positive definiteness of $B$. In particular,
\begin{equation}\label{eq:mu}
\mu(X)=\#H,\qquad \#H\mid\chi(X,L)\quad(L\in\Pic(X)).
\end{equation}
\end{theorem}
\begin{proof}
On $A$ every symmetric integral matrix class has a line bundle representative over $K$, constructed from the coordinate divisors and normalized Poincar\'{e} bundles. Under the product principal polarization, the homomorphism $\phi_L:A\to A^\vee$ of a representative is the matrix $B$. The theorem of the square identifies the effect of translation by $a$ on its Picard class with $\phi_L(a)$. Twisting the Galois action by the torsor cocycle shows that the component $\Pic^B(X)$ is an $A^\vee$-torsor of class $\pm B\xi$, where the sign depends on the translation convention. Hence that component has a rational point if and only if $B\xi=0$.

A rational Picard class has a possible Brauer obstruction to being represented by a line bundle. The low-degree Hochschild--Serre sequence puts the obstruction in $\ker(\Br(K)\to\Br(X))$. Evaluation at a point of $X(K_v)$ shows that it restricts to zero at every place $v$. Global Brauer reciprocity implies the obstruction is zero. Thus $B$ descends exactly when every row belongs to $\Lambda_\xi$, equivalently every column does, by symmetry. This proves \eqref{eq:pol}.

Over the algebraic closure, Riemann--Roch on an abelian variety identifies $\chi(X,L)$ with $c_1(L)^h/h!=\det B$. One can see the sign of the determinant by diagonalizing the symmetric matrix over $\QQ$ and using pullback by a matrix isogeny to reduce to a product of coordinate line bundles. The same diagonalization identifies ampleness with positive definiteness. If $\det B\ne0$, the index formula gives
\[
|\det B|=[\ZZ^h:B\ZZ^h]=\#H\,[\Lambda_\xi:B\ZZ^h].
\]
Divisibility holds also for zero determinant. Write the Smith normal form of the relation lattice as $\Lambda_\xi=UD\ZZ^h$, with $U\in\operatorname{GL}_h(\ZZ)$ and $D=\operatorname{diag}(q_1,\ldots,q_h)$, $q_i\mid q_{i+1}$. The matrix $UDU^{\mathsf T}$ is positive definite, has column lattice $\Lambda_\xi$, and has determinant $\prod_iq_i=\#H$. Its descended line bundle attains the minimum.
\end{proof}

\begin{corollary}[Stable obstruction]\label{cor:stable}
For every integer $t\geq0$, $\mu(X\times E^t)=\#H$. In particular a nontrivial torsor in this family cannot acquire a descended principal polarization line bundle by adding split elliptic factors.
\end{corollary}
\begin{proof}
The enlarged tuple consists of the original classes and $t$ zeros, and generates exactly $H$. Apply Theorem~\ref{thm:pol} to it.
\end{proof}

\begin{theorem}[Exact integral descent defect]\label{thm:def}
If $q_1\mid\cdots\mid q_h$ are the invariant factors of $H$ (including factors equal to one), then
\begin{equation}
\D(X/K):=\NS(X_{\overline K})/\NS_K(X)
\simeq\bigoplus_{j=1}^h(\ZZ/q_j\ZZ)^{\oplus j}.
\end{equation}
\end{theorem}
\begin{proof}
The congruence transformation $B\mapsto U^{-1}BU^{-\mathsf T}$ sends the condition in \eqref{eq:pol} to the condition that the columns lie in $D\ZZ^h$. The $(i,j)$ entry of a symmetric matrix must then be divisible by $\operatorname{lcm}(q_i,q_j)$. For $i\leq j$ this is $q_j$. Each of the $j$ independent entries in column $j$ in the upper triangle contributes one copy of $\ZZ/q_j\ZZ$ to the quotient.
\end{proof}

\section{Quantitative isogeny rigidity}
The polarization invariant does not directly control every morphism to an abelian variety, since the target might be isogenous to, but not isomorphic to, $A$. We now allow all targets of the same dimension. In this section the classes $\xi_i$ may belong to $H^1(K,E)$ without local solubility.
\begin{definition}
Let $\I(X)$ be the minimum degree of a dominant $K$-morphism $X\to B$, as $B$ ranges over all abelian varieties over $K$ of dimension $h$.
\end{definition}
Every such morphism is geometrically a translation of an isogeny $\varphi:A\to B$ defined over $K$, by the rigidity lemma. Equivariant torsor descent shows that such a morphism exists for $\varphi$ precisely when $\varphi_*\xi=0$ in $H^1(K,B)$. Thus the minimum is over these annihilating isogenies. For the fixed target $A$, the rows of the matrix of $\varphi$ must lie in $\Lambda_\xi$; a matrix with a row basis of $\Lambda_\xi$ attains determinant $\pm\#H$. Since the degree of a matrix isogeny is the square of its determinant, the fixed-target minimum is exactly $(\#H)^2$.

For a prime $p$, let $\mathcal A_p$ be the $\ZZ_p$-subalgebra of $M_2(\ZZ_p)$ spanned by the Galois image on $T_pE$. Serre's open image theorem \cite{Serre} gives a finite integer $c_p\geq0$ such that
\begin{equation}\label{eq:cp}
p^{c_p}M_2(\ZZ_p)\subseteq\mathcal A_p,
\end{equation}
with $c_p=0$ for all but finitely many primes. In fact openness gives a congruence subgroup, whose differences from the identity contain the required multiples of matrix units. If the residual image generates $M_2(\FF_p)$, Nakayama's lemma gives $c_p=0$.
\begin{lemma}[Matrix saturation]\label{lem:sat}
If $U\subset E[p^n]^h$ is a Galois-stable subgroup and $U^{\mathrm{sat}}=M_2(\ZZ/p^n\ZZ)U$, then
\[
p^{c_p}U^{\mathrm{sat}}\subseteq U,\qquad [U^{\mathrm{sat}}:U]\leq p^{2hc_p}.
\]
The group $U^{\mathrm{sat}}$ is the kernel of an integral matrix isogeny $A\to A$ of $p$-power degree.
\end{lemma}
\begin{proof}
Equation \eqref{eq:cp} and Galois stability imply the first inclusion. The quotient is killed by $p^{c_p}$ and requires at most $2h$ generators, because it is a quotient of a subgroup of $(\ZZ/p^n\ZZ)^{2h}$. This proves the index bound. Choose a basis of $E[p^n]$ and set $R=\ZZ/p^n\ZZ$. The diagonal and off-diagonal matrix units show that an $M_2(R)$-stable subgroup of $R^h\oplus R^h$ is $W\oplus W$ for some $W\subseteq R^h$. The annihilator lattice
\[
\Lambda_W=\{a\in\ZZ^h:a\cdot w=0\text{ in }R\text{ for all }w\in W\}
\]
has index $\#W$, by perfect finite-module duality. An integer matrix with row basis of $\Lambda_W$ has kernel $W\oplus W$ on $E[p^n]^h$. Its full geometric kernel is of order $(\#W)^2$, so this is the entire kernel. This gives the required isogeny.
\end{proof}

\begin{theorem}[Quantitative isogeny rigidity]\label{thm:iso}
There is a constant $C=C(E,K,h)=\prod_p p^{2hc_p}$, independent of $\xi$, such that
\begin{equation}\label{eq:iso}
\frac{(\#H)^2}{C}\leq\I(X)\leq(\#H)^2.
\end{equation}
If every prime dividing $\#H$ has $c_p=0$, then $\I(X)=(\#H)^2$.
\end{theorem}
\begin{proof}
The upper bound was shown for the fixed target. For the lower bound, let $\varphi:A\to B$ annihilate $\xi$, and let $U$ be the $p$-primary component of its kernel. The factorization through $A/U$ has remaining degree prime to $p$. An isogeny of degree prime to $p$ is injective on $p$-primary Weil--Ch\^{a}telet classes, since a complementary isogeny composes with it to give multiplication by an integer prime to $p$. Thus the quotient by $U$ annihilates the $p$-primary class $\xi_p$. The quotient by $U^{\mathrm{sat}}$ also annihilates it, and Lemma~\ref{lem:sat} identifies the latter quotient map, up to an isomorphism of targets, with a matrix isogeny $M_p$. Its rows annihilate the coordinate classes of $\xi_p$, which generate $H_p=H[p^\infty]$. Therefore
\[
(\#H_p)^2\leq(\det M_p)^2=\#U^{\mathrm{sat}}
\leq p^{2hc_p}\#U.
\]
Since $\#U$ is the $p$-part of $\deg\varphi$, multiplication over all primes proves the lower bound. For the exact assertion use only primes dividing $\#H$ in the lower bound, with $c_p=0$ for each. All other prime factors of $\deg\varphi$ are at least one.
\end{proof}
\begin{corollary}\label{cor:bound}
For fixed $E,K,h,D$, there are finitely many isomorphism classes of everywhere locally soluble $A$-torsors satisfying $\I(X)\leq D$. Their generated subgroups satisfy $\#H\leq\sqrt{CD}$.
\end{corollary}
\begin{proof}
By \eqref{eq:iso} the subgroup order is bounded. Every coordinate class is killed by $N!$, with $N=\lfloor\sqrt{CD}\rfloor$. For $N\geq1$ the group $\Sha(E/K)[N!]$ is finite by finite-level Selmer theory. Thus only finitely many tuples are possible. For $N=0$ there are none.
\end{proof}

\section{Canonical-height energy and algebraic curves}
Let $P\in E(K)$ have infinite order. Write $P_i$ for its embedding in the $i$-th coordinate of $A$. For a descent morphism $f:X\to B$ with associated isogeny $\varphi$ and a symmetric ample line bundle $L_B$, define the normalized canonical-height energy
\[
\mathcal E_P(f,L_B)=\frac{\sum_{i=1}^h\widehat h_{L_B}(\varphi P_i)}{\widehat h_{\mathcal O_E(0)}(P)}.
\]
This is an arithmetic intersection invariant expressed through canonical heights; it is not a bound on the stable Faltings height of the Albanese variety.
\begin{theorem}[Height energy]\label{thm:energy}
For any such $f,L_B$,
\[
\mathcal E_P(f,L_B)\geq h\left(\frac{(\#H)^2\chi(B,L_B)}{C}\right)^{1/h}.
\]
If the coordinate classes are independent of orders $q_i$, then for the fixed target $A$ with its product principal polarization,
\begin{equation}\label{eq:energy}
\min_f\mathcal E_P(f)=\sum_{i=1}^h q_i^2.
\end{equation}
\end{theorem}
\begin{proof}
The class of $\varphi^*L_B$ is a positive definite symmetric matrix $Q$. Functoriality of the canonical height and the coordinate embeddings give $\mathcal E_P=\operatorname{tr}Q$. By Riemann--Roch, $\det Q=(\deg\varphi)\chi(B,L_B)$. The arithmetic--geometric mean inequality on the eigenvalues of $Q$ and Theorem~\ref{thm:iso} give the lower bound. In the fixed-target case $Q=M^{\mathsf T}M$ for the descent matrix $M$. Independence forces the $i$-th column of $M$ to have entries divisible by $q_i$. Every column is nonzero, so its squared Euclidean norm is at least $q_i^2$. The diagonal matrix of the $q_i$ attains equality.
\end{proof}
Thus a bound $\mathcal E_P(f,L_B)\leq D$ implies $\#H\leq\sqrt C(D/h)^{h/2}$ and gives a second geometric finiteness statement by the argument of Corollary~\ref{cor:bound}.

\begin{theorem}[Genus spectrum]\label{thm:genus}
Let $X=X_1\times X_2$ be the product of independent everywhere locally soluble genus-one torsors of respective orders $q_1,q_2\geq3$. The genera of smooth geometrically integral curves embedded in $X$ over $K$ are exactly
\[
\{1+tq_1q_2:t=1,2,\ldots\}.
\]
\end{theorem}
\begin{proof}
The canonical bundle of $X$ is trivial. For a smooth curve $C\subset X$, adjunction and Theorem~\ref{thm:pol} give $g(C)-1=C^2/2=\chi(X,\mathcal O(C))$, divisible by $q_1q_2$. The genus is not zero, since there are no rational curves in an abelian variety in characteristic zero. If it were one, the two morphisms to the factors would give homomorphisms $J_C\to E$ and express both torsor classes as their pushforwards of $[\Pic^1(C)]$. The group $\Hom_K(J_C,E)$ has rank at most one because $E$ is geometrically non-CM. Both classes would lie in a common cyclic subgroup, contrary to independence.

For existence, Theorem~\ref{thm:pol} in dimension one gives line bundles $L_i$ on $X_i$ of degrees $q_i$. The bundle $L_1\boxtimes L_2^{\otimes t}$ is very ample over $K$ for every positive integer $t$. Bertini over an infinite field provides a smooth geometrically integral divisor in this system. Its genus is $1+tq_1q_2$ by adjunction.
\end{proof}
For $h\geq2$ independent coordinate classes of orders $q_i\geq3$, the external product of degree-$q_i$ line bundles is very ample and has Euler characteristic $\#H$. Intersecting $h-1$ general members produces a smooth complete-intersection curve of genus $1+(h-1)h!\#H/2$. These statements concern embedded curves or the indicated complete intersections, not all possible curves dominating the factors.

\section{Mixed Poincar\'{e} descent and integral motives}
Return to locally soluble $X$. On $X\times A$ the geometric N\'{e}ron--Severi group is given by symmetric matrices of size $2h$. Write a matrix in blocks as
\[
\begin{pmatrix}B&C\\ C^{\mathsf T}&D\end{pmatrix}.
\]
Let $\mathcal P_K(X,A)$ be the subgroup of mixed matrices $C$ occurring in line bundles on $X\times A$ defined over $K$.
\begin{theorem}[Poincar\'{e} obstruction]\label{thm:poincare}
There is an isomorphism $M_h(\ZZ)/\mathcal P_K(X,A)\simeq H^{\oplus h}$. The smallest positive multiple of the geometric mixed Poincar\'{e} class that descends is $\exp H$.
\end{theorem}
\begin{proof}
The class of the product torsor is $(\xi,0)$. By Theorem~\ref{thm:pol}, the descent conditions are $B\xi=0$ and $C^{\mathsf T}\xi=0$. The second says that each column of $C$ lies in $\Lambda_\xi$ and is sufficient by taking $B=D=0$. The quotient is therefore $h$ copies of $\ZZ^h/\Lambda_\xi$. The mixed Poincar\'{e} matrix is $I_h$ (up to the sign of the pairing); its multiple $mI_h$ descends exactly if $m\xi=0$.
\end{proof}
For two locally soluble genus-one torsors of classes $\xi,\eta$ under $E$, an analogous calculation on their product shows that the possible multipliers on the Jacobians induced by divisor correspondences form the ideal
\begin{equation}
\operatorname{lcm}\bigl(\ord(\xi\bmod\langle\eta\rangle),\ord(\eta\bmod\langle\xi\rangle)\bigr)\ZZ.
\end{equation}
Indeed an off-diagonal entry $m$ in a descending matrix exists precisely when $m\eta\in\langle\xi\rangle$ and $m\xi\in\langle\eta\rangle$.

Define the integral degree-two cycle defect by
\[
\Tdef_p(X/K)=\coker\bigl[\Pic(X)\otimes\ZZ_p\xrightarrow{c_1}
H^2_{\mathrm{et}}(X_{\overline K},\ZZ_p(1))^{G_K}\bigr].
\]
\begin{theorem}[Integral Tate-cycle comparison]\label{thm:tate}
There is a canonical comparison with the descended lattice,
\begin{equation}\label{eq:tate}
\Tdef_p(X/K)\simeq\D(X/K)\otimes\ZZ_p
\simeq\bigoplus_{j=1}^h(\ZZ_p/q_j\ZZ_p)^{\oplus j}.
\end{equation}
In particular, if $p\mid\#H$, the Chow motives of $X$ and $A$ with coefficients in $\ZZ_p$ are not isomorphic.
\end{theorem}
\begin{proof}
Translations act trivially on geometric \'etale cohomology. The twist by $\xi$ therefore does not change the Galois representation on degree two. K\"unneth gives $h$ coordinate degree-two components and one tensor-product component for each pair of distinct coordinates. The latter has invariants identified, via the perfect Weil pairing, with $\End_{G_K}(T_pE)=\ZZ_p$, by Faltings's endomorphism theorem \cite{Faltings}. The coordinate divisor classes and the Poincar\'{e} cross classes generate the resulting invariant lattice integrally: each coordinate class is primitive and the identity endomorphism generates the invariant cross tensor. Hence $\NS(X_{\overline K})\otimes\ZZ_p$ maps isomorphically onto this invariant lattice. The image of $\Pic(X)\otimes\ZZ_p$ is the descended lattice tensored with $\ZZ_p$, by flatness. Theorem~\ref{thm:def} gives the formula.

For $A$ the cycle defect vanishes. An isomorphism of integral Chow motives gives isomorphisms of codimension-one Chow groups and degree-two integral \'etale realizations, compatible with the cycle map by functoriality under correspondences. It would identify their defects, which is impossible when the right side of \eqref{eq:tate} is nonzero.
\end{proof}
Thus the integral descent data cannot be recovered from the geometric Galois representation alone. This statement does not claim any failure of rational Tate-cycle surjectivity.

\section{Selmer division fields: the Galois bridge}
In Sections 6--8 suppose that $p\geq5$ and $\rho_{E,p}(G_K)$ contains $\mathrm{SL}_2(\FF_p)$. Let $R_n=\ZZ/p^n\ZZ$, $V_n=E[p^n]$, $K_n=K(V_n)$, and $\Gamma_n=\Gal(K_n/K)$. Put $S_n=\Sel_{p^n}(E/K)$ and $M_n=E(K)/p^nE(K)$, viewed as a submodule of $S_n$. The Kummer sequence is
\begin{equation}\label{eq:kummer}
0\longrightarrow M_n\longrightarrow S_n\longrightarrow\Sha(E/K)[p^n]\longrightarrow0.
\end{equation}
For any finite submodule $D\subseteq H^1(K,V_n)$, define $L(D)/K_n$ by the common kernel of the restrictions of the classes in $D$ to $G_{K_n}$. These restrictions are well-defined homomorphisms into $V_n$, since the coefficient action is trivial over $K_n$. The field is finite Galois over $K$, and the evaluation group embeds in $\Hom_{R_n}(D,V_n)$.

\begin{lemma}[Full Selmer evaluation]\label{lem:eval}
There is a canonical evaluation isomorphism
\begin{equation}\label{eq:eval}
\Gal(L(D)/K_n)\simeq\Hom_{R_n}(D,V_n).
\end{equation}
The Galois action of $\Gamma_n$ on the right is its natural action on the target, with trivial action on $D$.
\end{lemma}
\begin{proof}
The residual image contains $-I$. Any lift of $-I$ to $\Gamma_n$ has order $2p^a$ for some $a\geq0$, since the kernel of reduction is a $p$-group. Its $p^a$-th power is an involution reducing to $-I$. The difference of that involution from $I$ is invertible, so the factorization $g^2-I=(g-I)(g+I)$ implies $g=-I$. The resulting central element acts as $-1$ on $V_n$. The standard central-element cohomological identity (Sah's lemma) yields $H^i(\Gamma_n,V_n)=0$ for all $i\geq0$, because $2$ is a unit. Restriction on $H^1(K,V_n)$ is therefore injective.

Write $U$ for the image of evaluation in $\Hom_{R_n}(D,V_n)$. Conjugation by $\Gamma_n$ acts on the values of the restricted cocycles, as follows immediately from the cocycle identity. The $R_n$-algebra generated by $\Gamma_n$ is $M_2(R_n)$: the elementary upper and lower unipotent matrices in the residual image generate the matrix units, and Nakayama's lemma lifts the assertion. Thus, after a basis choice for $V_n$, the module $U$ is $D'\oplus D'$ for a submodule $D'\subseteq\Hom_{R_n}(D,R_n)$. If $D'$ were proper, the perfect duality of finite $R_n$-modules would provide a nonzero $d\in D$ annihilated by all of $D'$. Its restricted cocycle would be zero, contradicting injectivity. Hence evaluation is surjective.
\end{proof}
\begin{theorem}[Relative Selmer division duality]\label{thm:bridge}
Let $H\subseteq\Sha(E/K)[p^n]$ be any subgroup, and let $S_{H,n}$ be its preimage under \eqref{eq:kummer}. Write $L_{H,n}=L(S_{H,n})$ and $L_{\mathrm{MW},n}=L(M_n)$. There is a canonical isomorphism
\begin{equation}\label{eq:bridge}
\Gal(L_{H,n}/L_{\mathrm{MW},n})\simeq\Hom_{R_n}(H,V_n).
\end{equation}
In particular $[L_{H,n}:L_{\mathrm{MW},n}]=(\#H)^2$. If $X$ is the product of any finite list of genus-one torsors generating $H$, then
\begin{equation}\label{eq:triple}
\boxed{[L_{H,n}:L_{\mathrm{MW},n}]=\I(X)=\mu(X)^2=(\#H)^2.}
\end{equation}
\end{theorem}
\begin{proof}
By Lemma~\ref{lem:eval}, the relative Galois group consists of exactly the evaluations on $S_{H,n}$ that vanish on $M_n$. By \eqref{eq:kummer}, this is $\Hom_{R_n}(H,V_n)$. Since $V_n$ is free of rank two over $R_n$, the cardinality is $(\#H)^2$. The polarization identity follows from Theorem~\ref{thm:pol}. The isogeny identity follows from Theorem~\ref{thm:iso}, since $c_p=0$ under the large-image hypothesis and $H$ is a $p$-group.
\end{proof}
The equality compares a Galois extension degree over the Mordell--Weil division field with the minimum degree of a morphism of varieties over $K$; it does not identify these two morphisms or their base fields. It holds also for nonhomocyclic $H$. When $H=(\ZZ/p^n)^h$, the three degrees in \eqref{eq:triple} all equal $p^{2nh}$.

\begin{corollary}[Affine capitulation]\label{cor:cap}
If $H\simeq R_n^h\subseteq\Sha(E/K)[p^n]$, there exists a finite extension $F/K$ of degree $(\#H)^2=p^{2nh}$ annihilating every class of $H$. It is unramified outside the primes above $p$ and the primes of bad reduction. No minimality of this field degree is asserted.
\end{corollary}
\begin{proof}
Choose lifts $c_1,\ldots,c_h\in S_n$ of a basis of $H$, represented by cocycles. Their classes generate a free $R_n$-module of rank $h$ by independence of the images in $H$. The affine representation $g\mapsto((c_i(g))_i,\rho_{E,p^n}(g))$ has full image $V_n^h\rtimes\Gamma_n$: the translation image is all of $V_n^h$ by Lemma~\ref{lem:eval}, and the linear image is surjective. In its Galois field take the fixed field of the subgroup of zero translations, isomorphic to $\Gamma_n$. The index is $p^{2nh}$, and the chosen cocycles restrict to zero there. At good primes away from $p$ the local Selmer condition is unramified, as is the torsion representation; every cocycle vanishes on inertia. The ramification assertion follows. The field $F$ is in general not Galois over $K$.
\end{proof}

\section{Infinite Selmer duality and Lie algebras}
Let $K_\infty=K(E[p^\infty])$, and define $L_\infty^{\mathrm{Sel}}$ and $L_\infty^{\mathrm{MW}}$ by adjoining $K_\infty$ to the composita of $L(S_n)$ and $L(M_n)$, respectively. The compatible scalar $-I$ lies in $\Gal(K_\infty/K)$ by compactness, so restriction from $K$ to $K_\infty$ is injective in degree one at every finite precision. The same proof as Lemma~\ref{lem:eval} gives full evaluation over $K_\infty$.

Write $\mathcal S=\Sel_{p^\infty}(E/K)$, $\mathcal M=E(K)\otimes\QQ_p/\ZZ_p$, and $\mathcal A=\Sha(E/K)[p^\infty]$. The condition $E(K)[p]=0$ implies that the finite-level classes embed compatibly into $\mathcal S$ and form its $p^n$-torsion. Passing to inverse limits of the evaluation groups and using the annihilator of $\mathcal M$ gives the following result.
\begin{theorem}[Infinite division duality]\label{thm:infty}
With the pointwise convergence topology on Hom, there is a canonical isomorphism
\begin{equation}
\mathfrak X_p:=\Gal(L_\infty^{\mathrm{Sel}}/L_\infty^{\mathrm{MW}})
\simeq\Hom(\mathcal A,E[p^\infty]).
\end{equation}
If $d_p=\corank_{\ZZ_p}\mathcal A$, then $\rk_{\ZZ_p}\mathfrak X_p=2d_p$.
\end{theorem}
\begin{proof}
At each finite level, full evaluation gives all homomorphisms from $S_n$ to $V_n$. The transition maps are restriction of homomorphisms under $S_n\hookrightarrow S_{n+1}$ and $V_n\hookrightarrow V_{n+1}$, since the restricted cocycles agree under coefficient inclusion. Finite-module duality (or injectivity of $\QQ_p/\ZZ_p$) ensures extension of each homomorphism. The inverse limit is consequently $\Hom(\mathcal S,E[p^\infty])$. The condition of fixing all Mordell--Weil division fields is exactly annihilation of $\mathcal M$. The infinite Kummer sequence identifies the quotient with $\mathcal A$. Finally the cofinitely generated group $\mathcal A$ is a direct sum of $(\QQ_p/\ZZ_p)^{d_p}$ and a finite group, giving the rank formula.
\end{proof}
At every finite place $v\nmid p$, the local group $E(K_v)$ is an extension of a finite group by an open subgroup uniquely divisible by $p$, so $E(K_v)\otimes\QQ_p/\ZZ_p=0$. The infinite Selmer condition therefore implies that $L_\infty^{\mathrm{Sel}}/K_\infty$ splits completely at all such places, even at bad primes. At primes above $p$ the additional strict local condition is necessary for complete splitting.

Let $V=T_pE\otimes\QQ_p$ and write $\mathcal F$ for the rational fine Selmer space, the kernel of localization at primes above $p$ on the rational classical Selmer group. Write $\mathcal P$ for the kernel of the corresponding localization on $E(K)\otimes\QQ_p$. The fine quotient need not have dimension $d_p$ over an arbitrary number field: there can be an additional local quotient. Over $\QQ$ of positive rank $r$, however, the local Kummer space is one-dimensional and the elliptic logarithm of any non-torsion rational point is nonzero. Thus $\dim\mathcal P=r-1$ and $\dim\mathcal F=r-1+d_p$.
\begin{theorem}[Affine Selmer Lie algebra]\label{thm:lie}
For a basis $c_1,\ldots,c_s$ of $\mathcal F$ represented by rational cocycles, let $\widetilde G$ be the image of the associated affine representation into $\operatorname{GL}(V)\ltimes V^s$. Then
\[
\operatorname{Lie}(\widetilde G)=\mathfrak{gl}(V)\ltimes\Hom_{\QQ_p}(\mathcal F,V).
\]
The relative translation Lie algebra annihilating $\mathcal P$ is $\Hom_{\QQ_p}(\mathcal F/\mathcal P,V)$. Over $\QQ$ in positive rank its dimension is $2d_p$.
\end{theorem}
\begin{proof}
Serre's theorem gives the full reductive Lie algebra $\mathfrak{gl}(V)$. Over $K_\infty$, take the rational span of the joint restricted cocycle image. It is invariant under the full matrix algebra on $V$. As in Lemma~\ref{lem:eval}, a proper subspace would annihilate a nonzero class in $\mathcal F$, contrary to inflation--restriction and scalar vanishing. The compact translation image has Lie algebra equal to its rational span. Thus the full translation ideal occurs and the projection onto the linear Lie algebra is surjective, proving the formula. Annihilating $\mathcal P$ gives the relative assertion.
\end{proof}
For the Lie algebra $\mathfrak g=\mathfrak{gl}(V)\ltimes\Hom(\mathcal F,V)$ with translation ideal acting trivially on $V$, Hochschild--Serre and the central identity matrix give $H^1_{\mathrm{Lie}}(\mathfrak g,V)\simeq\mathcal F$. Indeed the cohomology of $\mathfrak{gl}(V)$ with coefficients in $V$ vanishes by the Cartan homotopy formula, leaving $\Hom_{\mathfrak{gl}(V)}(V\otimes\mathcal F^*,V)\simeq\mathcal F$. This does not force the translation ideal to vanish: semidirect Lie algebras with any number of additional copies of $V$ have the same reductive component.

\section{Fine Selmer class-field comparison}
For the following exact finite-level comparison impose $K=\QQ$, positive rank $r$, good ordinary non-anomalous reduction at $p$, and $E(\QQ_\ell)[p]=0$ at bad primes $\ell\ne p$. Let $S$ contain $p$, the bad primes and the archimedean place; let $C_n=\Cl_S(K_n)$ be the ideal class group modulo primes above $S$. Define
\[
h_n=\length_{\ZZ_p}\Hom_{\Gamma_n}(C_n,V_n).
\]
\begin{theorem}[Exact strict multiplicity]\label{thm:fine}
Let $a$ denote the index exponent of the image of $E(\QQ)\otimes\ZZ_p$ in the local $p$-adic completion $\ZZ_p$ of $E(\QQ_p)$. Then
\begin{equation}\label{eq:hn}
h_n=n(r-1)+\log_p\#\Sha(E/\QQ)[p^n]+b_n,
\end{equation}
where $0\leq b_n\leq\min(a,n)$ and $b_n$ is nondecreasing. In particular $\lim h_n/n=r-1+d_p$. If $a=0$ then $b_n=0$ identically.
\end{theorem}
\begin{proof}
At $p$, the ordinary residual representation is triangular with diagonal characters $\omega\psi^{-1},\psi$, where $\omega$ is cyclotomic and $\psi$ unramified. Let $\beta=\psi(\mathrm{Frob}_p)\in\FF_p^\times$, so non-anomalous means $\beta\ne1$. Multiplying a Frobenius lift by suitable inertia yields a residual matrix with both diagonal entries $\beta$; its $p$-th power is $\beta I$. A nonidentity residual scalar lifts to an actual central scalar in each finite local image: take the prime-to-$p$ part of its order, then apply the factorization of $X^m-1$ at the Teichm\"uller lift of $\beta$, with remaining factor invertible modulo $p$. The central scalar minus one acts invertibly, giving zero local inflation groups. At bad primes away from $p$ the full local $H^1(\QQ_\ell,V_n)$ vanishes by local duality and Euler characteristic, using the assumed absence of local $p$-torsion. The real place contributes nothing since $p$ is odd.

Global inflation and transgression vanish by the central $-I$ argument. Class-field theory identifies the strict finite-level Selmer group $R_n^0$ (all local conditions at $S$ zero, unramified outside $S$) exactly with $\Hom_{\Gamma_n}(C_n,V_n)$: restriction is injective; descent follows from zero transgression, and the local inflation vanishing makes the descended class strict. This is the split-at-$S$ version of the representation-specific class-field comparison; compare \cite{HO} for related Iwasawa-theoretic results.

The full Selmer group has length $nr+\log_p\#\Sha[p^n]$. Its local Kummer image at $p$ has order $p^{n-b_n}$ for some $b_n$, and the other local conditions at $S$ are zero. The local Mordell--Weil image has index exponent $\min(a,n)$ at precision $n$, so $b_n$ is no larger. Reduction from precision $n+1$ to $n$ sends local Selmer images into each other. It follows that $b_{n+1}\geq b_n$ (if $b_{n+1}>n$ this is automatic; otherwise use the image of $p^{b_{n+1}}$). Subtracting the local image length proves \eqref{eq:hn}. The limit follows from the structure of cofinitely generated $p$-primary groups.
\end{proof}
In particular the difference $h_{n+1}-h_n$ is at least $r-1+d_p$. If it equals $r-1$ at any finite precision, then $d_p=0$ and $p^n\Sha(E/\QQ)[p^\infty]=0$: writing the finite summands with exponents $e_i$, the increment minus $r-1$ is $d_p+\#\{i:e_i>n\}+b_{n+1}-b_n$, a sum of nonnegative integers. This is a finite arithmetic certificate, not a universal bound on the increment.

\section{Compatible torsor towers and Brauer obstructions}
This section assumes only the geometrically non-CM hypothesis and $p\geq5$. Suppose that the divisible part of $\Sha(E/K)[p^\infty]$ contains $(\QQ_p/\ZZ_p)^h$. Choose coordinate classes $\xi_{i,n}$ of order $p^n$ with $p\xi_{i,n+1}=\xi_{i,n}$ and independent at every level. Write $X_n$ for the product torsor attached to these classes, and $X_0=A$. Equivariant pushforward along $[p]$ gives morphisms $X_{n+1}\to X_n$ which are finite \'etale of degree $p^{2h}$. The geometric N\'{e}ron--Severi lattice is constant, but its descended sublattice is $p^n\Sym_h(\ZZ)$ by Theorem~\ref{thm:pol}.

For a fixed $X=X_n$ and any $m\geq1$, the composite $X_{n+m}\to X_n$ is an $A[p^m]$-torsor. Its class $u_m$ in $H^1_{\mathrm{et}}(X,A[p^m])$ is compatible with coefficient multiplication by $p$ in the tower. Write $u_{i,m}$ for its $i$-th coordinate. The Weil pairing produces classes
\[
\theta_{ij,m}=e_{p^m,*}(u_{i,m}\smile u_{j,m})\in H^2_{\mathrm{et}}(X,\mu_{p^m}).
\]
Geometrically the coverings are multiplication coverings, so the diagonal classes $\theta_{ii,m}/2$ and the off-diagonal classes $\theta_{ij,m}$ for $i<j$ restrict to the coordinate and mixed divisor classes, with signs depending on the Weil pairing convention. This follows by evaluating the identity class in $\Hom(E[p^m],E[p^m])$ on the geometric fundamental group and taking its alternating pairing; on the diagonal the cup product contributes twice the standard generator. These classes are compatible under coefficient reduction because the Weil pairings are compatible with the transition maps. Consequently they lift a basis of the invariant integral Tate lattice to the inverse limit of the arithmetic groups $H^2_{\mathrm{et}}(X,\mu_{p^m})$. In contrast, the image of actual line bundles is $p^n\Sym_h(\ZZ_p)$. This distinction is related to the existence of abelianized fundamental group sections for divisible Albanese classes \cite{HS}.

Define the locally trivial algebraic Brauer group by
\[
\mathcal B(X)=\ker\left[\Br_1(X)\to\prod_v\Br(X_{K_v})\right],\qquad
\Br_1(X)=\ker[\Br(X)\to\Br(X_{\overline K})].
\]
Here $T_p$ of a torsion group is its inverse limit of $p^m$-torsion under multiplication by $p$.
\begin{theorem}[Brauer lifting of the polarization defect]\label{thm:brauer}
For the compatible tower just constructed and every $n\geq1$, there is a surjection
\begin{equation}\label{eq:brauer}
T_p\mathcal B(X_n)\twoheadrightarrow\Tdef_p(X_n/K)
\simeq(\ZZ/p^n\ZZ)^{h(h+1)/2}.
\end{equation}
\end{theorem}
\begin{proof}
The Kummer sequence yields
\[
0\to\Pic(X)/p^m\Pic(X)\to H^2_{\mathrm{et}}(X,\mu_{p^m})
\to\Br(X)[p^m]\to0.
\]
The constructed cup products have geometric restrictions in the image of the geometric Picard group, so their Brauer images are algebraic. They are locally zero in the Brauer group: at a place $v$, the nonempty compact sets $X_j(K_v)$ have a compatible sequence of origins by the finite intersection property. Under these origins the coverings become ordinary multiplication coverings. The cup products are then eigenclasses of $[2]^*$ with eigenvalue $4$; the corresponding symmetric coordinate and Poincar\'{e} line bundles have cycle classes with the same eigenvalue. Their difference lies in the first Hochschild--Serre filtration of arithmetic degree-two cohomology. On its two possibly nonzero graded pieces the eigenvalues are $2$ and $1$, respectively, because $[2]$ acts as $2^j$ on geometric degree $j$. Since $4-2$ and $4-1$ are units modulo $p^m$, the difference vanishes, by applying the filtration successively. Hence the local cup products are cycle classes of line bundles, and their Brauer images are zero.

For an algebraic Brauer class of order dividing $p^m$, lift it to arithmetic Kummer cohomology. Its geometric restriction lies in $\Pic(X_{\overline K})/p^m=\NS(X_{\overline K})/p^m$ because the geometric Picard variety is divisible. Modulo $\NS_K(X)$ and $p^m$, that restriction is independent of the lift. This construction gives compatible maps
\[
\mathcal B(X)[p^m]\to
\frac{\NS(X_{\overline K})}{\NS_K(X)+p^m\NS(X_{\overline K})}.
\]
The cup-product classes yield compatible Brauer classes whose images span the inverse limit of these quotients. The inverse limit is the defect in \eqref{eq:tate}. This gives the surjection.
\end{proof}
\begin{theorem}[Locally trivial Brauer corank]\label{thm:brank}
For every locally soluble torsor $X$ under $E^h$, without a tower assumption,
\begin{equation}\label{eq:brank}
\rk_{\ZZ_p}T_p\mathcal B(X)=h\,d_p,
\qquad d_p=\corank_{\ZZ_p}\Sha(E/K)[p^\infty].
\end{equation}
\end{theorem}
\begin{proof}
The geometric Picard sequence has connected part $A^\vee$ and constant torsion-free quotient $\NS(X_{\overline K})$. Since $H^1(K,\NS(X_{\overline K}))=0$, it gives
\[
H^1(K,\Pic(X_{\overline K}))\simeq H^1(K,A^\vee)/D_X,
\]
where $D_X$ is the finite image of the connecting homomorphism, contained in $\Sha(A^\vee/K)$ by local solubility. Locally the connecting map is zero because a local point identifies the torsor with $A$ and its matrix divisor classes descend. The Hochschild--Serre sequence, using $H^3(K,\overline K^\times)=0$, identifies $\Br_1(X)/\Br(K)$ with this Picard cohomology group; constants inject by local solubility and global reciprocity. It follows that the classes which restrict locally to constants correspond exactly to $\Sha(A^\vee/K)/D_X$. For such a class the sum of local invariants of a representative is independent of the representative. The kernel of this sum is precisely $\mathcal B(X)$: when the sum is zero the local constants come from a global constant by Brauer--Hasse--Noether, and subtraction gives a locally trivial representative. The local invariants have finite support by properness and unramifiedness outside a finite set. A global zero-cycle of degree $I>0$ shows that the image of the sum is killed by $I$, by evaluation and reciprocity, so is finite. These are the algebraic Brauer comparisons underlying the Albanese--Picard pairing (cf. \cite{PS}). Hence taking $p$-primary coranks gives $\corank\mathcal B(X)[p^\infty]=\corank\Sha(A^\vee/K)[p^\infty]=hd_p$. These groups are cofinitely generated; the corank is the rank of the Tate module.
\end{proof}
Since the target of \eqref{eq:brauer} requires $h(h+1)/2$ generators, \eqref{eq:brank} implies $h(h+1)/2\leq hd_p$. This is not a contradiction: the construction required $h\leq d_p$. Even the quadratic amplification of the divisor obstruction is consistent with a positive divisible corank.

\section{Kato multiplicities and the remaining reciprocity problem}
Over $\QQ$ with positive Mordell--Weil rank $r$, put $s_p=\corank R(E/\QQ)$ for the fine $p^\infty$-Selmer group (strict at $p$). The local infinite Kummer condition at every prime away from $p$ vanishes. At $p$ the local Kummer module has corank one, and the image of the global Mordell--Weil module has corank one by the nonzero elliptic logarithm of a non-torsion point. The infinite Kummer exact sequence therefore gives $s_p=r-1+d_p$. Under the hypotheses of Theorem C and Corollary C of \cite{BCGS} (non-CM, odd good ordinary $p$, no local $p$-torsion at $p$, and the rational cyclotomic main conjecture, including their established cases), the minimum nonzero Kolyvagin derivative order equals $s_p$. It follows from Theorem~\ref{thm:brank} that
\begin{equation}\label{eq:kato}
\ord(\boldsymbol\kappa^{\mathrm{Kato}})=r-1+d_p
=r-1+\frac{1}{h}\rk_{\ZZ_p}T_p\mathcal B(X).
\end{equation}
Here the order is the minimum number of auxiliary prime factors indexing a nonzero derivative at some permitted precision. Under the hypotheses of Kim's Kurihara-number structure theorem \cite{Kim}, the corresponding minimum modular-symbol order is $r+d_p$. These are numerical consequences of existing Euler-system structure theorems, not an identification of individual motivic classes with divisor cycles. The hypotheses and conventions of the cited theorems must hold before these consequences can be used.

In the exact local setting of Theorem~\ref{thm:fine}, we additionally have the identity
\[
\ord(\boldsymbol\kappa^{\mathrm{Kato}})=\lim_{n\to\infty}\frac{h_n}{n},
\]
whenever the Kato order formula is applicable. Thus the class-field growth, relative Selmer Galois rank and locally trivial Brauer rank all detect the same extra arithmetic multiplicity. Higher explicit reciprocity would require a compatibility on the level of motivic zeta elements and arithmetic regulators, rather than the dimension identity \eqref{eq:kato}. Burns, Kurihara and Sano's generalized Perrin--Riou conjecture \cite{BKS} provides a precise setting for this much deeper problem.

\section{Unified growth laws and finiteness criteria}
Let $H_n=\Sha(E/K)[p^n]$. Choose a product of genus-one torsors $X_n'$ representing generators of $H_n$. The number of generators can be bounded independently of $n$ by the finite-level structure of the cofinitely generated group. We can use a fixed number $h$ of factors by padding the generating list with zero classes. Put $\mu_n=\mu(X_n')$ and $D_n=[L_{H_n,n}:L_{\mathrm{MW},n}]$ at a large-image prime. If $\Sha(E/K)[p^\infty]\simeq(\QQ_p/\ZZ_p)^d\oplus\bigoplus_i\ZZ/p^{a_i}\ZZ$, the preceding results give
\begin{align}
\log_p\mu_n&=nd+\sum_i\min(n,a_i),\\
\log_p\I(X_n')&=2nd+2\sum_i\min(n,a_i)+O(1),\\
\log_p D_n&=2nd+2\sum_i\min(n,a_i)\qquad\text{(large image)}.
\end{align}
The error in the second formula is bounded by $\log_p C(E,K,h)$ and disappears at large-image primes. In particular the slope of the relative Galois degree is exactly twice the slope of the polarization invariant. If a divisible tower with $h=d>0$ exists, the exact cycle defect has order $p^{nd(d+1)/2}$, exhibiting a quadratic amplification in the divisible corank.

\begin{center}
\setlength{\fboxsep}{10pt}\fcolorbox{teal}{white}{\begin{minipage}{.87\textwidth}\centering
\textbf{The arithmetic--geometric--Galois bridge}\\[6pt]
$\displaystyle H\subseteq\Sha(E/K)[p^n]$\\[4pt]
$\Downarrow$\\[3pt]
$\displaystyle\mu(X)=\#H\qquad\quad\Gal(L_{H,n}/L_{\mathrm{MW},n})=\Hom(H,E[p^n])$\\[6pt]
$\displaystyle\I(X)=\mu(X)^2=[L_{H,n}:L_{\mathrm{MW},n}]$\\[6pt]
{\small Exact at large-image primes; the Galois group equality denotes the canonical evaluation isomorphism.}
\end{minipage}}
\end{center}

\begin{corollary}[Equivalent boundedness problems]\label{cor:equiv}
For fixed $E,K,h\geq1$, the following statements are equivalent:
\begin{enumerate}
\item $\Sha(E/K)$ is finite.
\item $\mu(X)$ is uniformly bounded over everywhere locally soluble $E^h$-torsors.
\item $\I(X)$ is uniformly bounded over the same family.
\end{enumerate}
For a fixed large-image prime $p$, the analogous equivalence holds for the $p$-primary subgroup and for the boundedness of the relative Selmer division-field degrees at all precisions.
\end{corollary}
\begin{proof}
If $\Sha$ is finite then $\#H$ is bounded, so the bounds follow. Conversely, a uniform bound on $\mu$ or $\I$ bounds the orders of all individual Sha classes by applying the invariant to a tuple with only one nonzero coordinate. A fixed factorial then kills all of them, so finite-level Selmer theory gives finiteness. For the $p$-primary assertion use the same argument and \eqref{eq:bridge}.
\end{proof}
These equivalences do not supply independent upper bounds. For example, all the torsors have the same geometric underlying abelian variety, the same geometric cohomology, and the same stable Faltings height of their Albanese variety. Faltings's finiteness theorem for abelian varieties does not count torsors under a fixed abelian variety. A uniform degree or integral cycle bound would require additional arithmetic input.

\section{Conclusion}
The central contribution is the quantitative comparison of geometric and Galois degrees, including exact equality for large-image primes. The seven principal geometric results are established with their stated hypotheses, and the earlier Galois bridge is made explicit in Theorems~\ref{thm:bridge}--\ref{thm:lie}. The tower calculation isolates a locally trivial Brauer obstruction even when all integral geometric degree-two classes lift to arithmetic \'etale cohomology. The Kato comparison is limited to numerical multiplicities obtained from established structural results. A uniform bound on the integral descent defects, or an independent higher motivic reciprocity forcing rank-sensitive nonvanishing, remains open. No numerical division-field class-group calculation or new Euler-system nonvanishing theorem is asserted. The derivations are intended for independent mathematical review; historical priority for their combination and consequences has not been established.

\clearpage
\begin{thebibliography}{99}\small
\bibitem{BCGS} A. Burungale, F. Castella, G. Grossi and C. Skinner, \emph{Non-vanishing of Kolyvagin systems and Iwasawa theory}, arXiv:2312.09301, \href{https://arxiv.org/abs/2312.09301}{revised version (2026)}.
\bibitem{BKS} D. Burns, M. Kurihara and T. Sano, \emph{On derivatives of Kato's Euler system for elliptic curves}, J. Math. Soc. Japan \textbf{76} (2024), 855--884, \href{https://doi.org/10.2969/jmsj/90699069}{doi:10.2969/jmsj/90699069}.
\bibitem{Faltings} G. Faltings, \emph{Endlichkeitss\"atze f\"ur abelsche Variet\"aten \"uber Zahlk\"orpern}, Invent. Math. \textbf{73} (1983), 349--366, \href{https://doi.org/10.1007/BF01388432}{doi:10.1007/BF01388432}.
\bibitem{HS} D. Harari and T. Szamuely, \emph{Galois sections for abelianized fundamental groups}, Math. Ann. \textbf{344} (2009), 779--800, \href{https://arxiv.org/abs/0808.2556}{arXiv:0808.2556}.
\bibitem{HO} T. Hiranouchi and T. Ohshita, \emph{Asymptotic behavior of class groups and cyclotomic Iwasawa theory of elliptic curves}, J. Th\'eor. Nombres Bordeaux (2023), \href{https://doi.org/10.5802/jtnb.1258}{doi:10.5802/jtnb.1258}.
\bibitem{Kim} C.-H. Kim, \emph{The structure of Selmer groups and the Iwasawa main conjecture for elliptic curves}, \href{https://arxiv.org/abs/2203.12159}{arXiv:2203.12159}.
\bibitem{MR} B. Mazur and K. Rubin, \emph{Kolyvagin Systems}, Memoirs Amer. Math. Soc. \textbf{168}, no. 799 (2004).
\bibitem{Milne} J. S. Milne, \emph{Arithmetic Duality Theorems}, 2nd ed., BookSurge, 2006.
\bibitem{Mumford} D. Mumford, \emph{Abelian Varieties}, Oxford University Press, 1970.
\bibitem{PS} B. Poonen and M. Stoll, \emph{The Cassels--Tate pairing on polarized abelian varieties}, Ann. of Math. \textbf{150} (1999), 1109--1149, \href{https://arxiv.org/abs/math/9911267}{arXiv:math/9911267}.
\bibitem{RS} J. Rosen and A. Shnidman, \emph{N\'{e}ron--Severi groups of product abelian surfaces}, \href{https://arxiv.org/abs/1402.2233}{arXiv:1402.2233}.
\bibitem{Serre} J.-P. Serre, \emph{Propri\'et\'es galoisiennes des points d'ordre fini des courbes elliptiques}, Invent. Math. \textbf{15} (1972), 259--331, \href{https://doi.org/10.1007/BF01405086}{doi:10.1007/BF01405086}.
\end{thebibliography}
\end{document}
